\documentclass[12pt]{article}
\usepackage[utf8]{inputenc}
\usepackage[T1]{fontenc}
\usepackage{hyperref}
\title{Karankawa Language Resources for K-12 Curriculum Development}
\author{Compiled by Claude Sonnet 5.5 (claude-sonnet-5-5) for Andr\'as Kornai}
\date{}
\begin{document}
\maketitle

\section{Introduction}
Karankawa (ISO 639-3 zkk) is an extinct language isolate of the Texas Gulf coast. Only a few hundred words survive. The language is being revitalized by community groups, but no school curriculum was found in the searches.

\section{Primary Educational Resources}
No primers, textbooks or school programmes were found. The Karankawa Tribe of Texas has a Language Keeper who teaches classes to tribal members and reconnecting relatives. The page is at \url{https://karankawas.com/language/}.

\section{Linguistic and Reference Materials}
\begin{itemize}
\item Wikipedia overview of the Karankawa language: \url{https://en.wikipedia.org/wiki/Karankawa_language}
\item Glottolog record for Albert S. Gatschet's 1891 work, a key source: \url{https://glottolog.org/resource/reference/id/307940}
\item Smithsonian finding aid for a manuscript collection (NAA MS 1014B): \url{https://sirismm.si.edu/EADpdfs/NAA.MS1014B.pdf}
\item EBSCO research starter on the Karankawa language: \url{https://ebsco.com/research-starters/language-and-linguistics/karankawa-language}
\end{itemize}

\section{Religious and Community Resources}
\begin{itemize}
\item Karankawa Tribe of Texas, language page: \url{https://karankawas.com/language/}
\item Karankawa Tribe of Texas, frequently asked questions: \url{https://karankawas.com/frequently-asked-questions/}
\end{itemize}

\section{Digital and Audio Resources}
No audio, video or software resources were found.

\section{Curriculum Development Context}
The word lists are mostly secondary and incomplete. The sources include Gatschet (1891), Swanton (1940) and earlier French lists, and accounts differ on who supplied Gatschet's material. A heritage curriculum would need to be built with the tribal Language Keepers, starting from a small core vocabulary. Check original sources before using any entry.
\end{document}